\chapter{Derivatives Markets and Contingent Claims}
\label{chap:Derivatives_Markets_and_Contingent_Claims}

\begin{openingcase}{How Do Derivatives Markets Organize Risk Transfer?}

A manufacturer with operations in several countries uses foreign-exchange forwards, interest-rate swaps, and commodity options to manage different financial exposures. Some contracts are standardized and traded on organized exchanges, where central clearing, margining, and daily settlement support market-wide trading. Others are negotiated bilaterally with dealers to match the firm's currency, maturity, notional amount, and payment dates more closely.

The firm therefore uses more than one type of derivatives market. The choice affects contract flexibility, price transparency, liquidity, counterparty exposure, collateral requirements, and the way positions are valued and settled. A customized contract may provide a closer hedge, while a standardized contract may be easier to trade or clear.

The underlying business exposures remain with the firm. What changes is the market architecture through which selected risks are transferred, priced, and managed.

\end{openingcase}



\begin{tcolorbox}[
  enhanced,
  colback=FinanceBlue!6,
  colframe=FinanceBlue,
  boxrule=0.6pt,
  sharp corners,
  left=8pt,
  right=8pt,
  top=7pt,
  bottom=7pt,
  before skip=1em,
  after skip=1.6em,
  title=\textbf{Core Questions},
  coltitle=white,
  colbacktitle=FinanceBlue,
  fonttitle=\bfseries,
  attach boxed title to top left={xshift=6pt,yshift=-2mm},
  boxed title style={
    boxrule=0pt,
    sharp corners,
    left=5pt,
    right=5pt,
    top=2pt,
    bottom=2pt
  }
]
\begin{enumerate}[
  leftmargin=*,
  label=\textcolor{FinanceBlue}{\textbf{\arabic*.}},
  itemsep=0.55em,
  topsep=0.2em
]


\item How do derivatives use state-contingent claims to unbundle and reallocate financial exposures?

\item How do no-arbitrage and replication connect derivative prices to the prices of other traded claims?

\item How do exchange trading, bilateral contracting, collateral, netting, and central clearing affect counterparty risk and the allocation of risk in derivative markets?


\end{enumerate}
\end{tcolorbox}

\section[Why Derivatives Exist: Contingent Claims and Risk Allocation]{Why Derivatives Exist:\\ Contingent Claims and Risk Allocation}

A derivative is a financial contract whose payoff depends on the value of an underlying asset, rate, index, event, or other specified condition. Its economic role is not confined to speculation. Derivatives allow participants to separate exposures that are bundled together in ordinary business and financial positions and to transfer selected parts of those exposures to other market participants.

\subsection{State-Contingent Payoffs}

A contingent claim is a financial claim whose payoff
depends on the future state of the world. Let the
possible states at date $T$ be indexed by
$s=1,\ldots,S$. The payoff of a claim in state $s$
is denoted by $X_T(s)$.

A state represents the future conditions relevant
to the contractual payoff. These conditions may
include the price of an underlying asset, an interest
rate, an exchange rate, a commodity price, or the
occurrence of a specified credit event. A contingent
claim therefore specifies how a contractual payment
varies across possible future states.

A forward contract illustrates a linear
state-contingent payoff. It specifies the purchase
and sale of an underlying asset at a future date $T$
for a delivery price $K$ agreed upon at inception.
The party obligated to purchase the underlying asset
holds the long position, while the party obligated
to deliver it holds the short position.

At maturity, the payoff per unit of the underlying
asset to the long forward position is

\begin{equation}
X_T^{\mathrm{long}}
=
S_T-K.
\label{eq:ch11-forward-long-payoff}
\end{equation}

The short forward position has the opposite payoff:

\begin{equation}
X_T^{\mathrm{short}}
=
K-S_T
=
-X_T^{\mathrm{long}}.
\label{eq:ch11-forward-short-payoff}
\end{equation}

Here, $S_T$ denotes the market price of the underlying
asset at maturity. The two positions have equal and
opposite payoffs before transaction costs and
counterparty losses.

Options also create state-contingent claims, but
their holders possess exercise rights rather than
an unconditional obligation to transact. Their
nonlinear payoff structures are examined in
later chapters.

The same state-contingent perspective applies more
broadly to debt, equity, insurance contracts, and
guarantees. These claims distribute payments
differently across future states, reflecting
differences in contractual priority, risk allocation,
and the conditions governing payment.

\subsection{Exposure Decomposition and Risk Unbundling}

Financial and operating positions often contain several sources of risk at the same time, and these exposures may interact. For example, the cash flow of an importing firm may be affected by both commodity prices and exchange rates, while joint movements in the two variables can amplify or partially offset changes in the firm's domestic-currency input costs. The risk of a position therefore depends not only on its individual exposures, but also on how those exposures move together.

Exposure decomposition separates an overall position into economically distinct sources of risk and identifies how each source affects the position's cash flows or value. This makes it possible to distinguish exposures that may be retained from those that may be reduced or otherwise modified.

Once distinct exposures have been identified, derivatives can be used to modify them separately. A derivative linked to a particular price, rate, index, or event creates a payoff that responds to that source of risk without requiring the underlying business activity or asset position to be changed. In this way, an exposure embedded in a broader position can be unbundled and modified separately from the other components of the position.

To illustrate risk unbundling, suppose an existing position has value $V_T$ at date $T$ and is exposed to several risk factors. A derivative linked to one of these factors has payoff $X_T$. The combined position is

\begin{equation}
V_T^{D}=V_T+X_T.
\label{eq:ch11-combined-derivative-position}
\end{equation}

Because $X_T$ responds to a selected risk factor, that component of the original exposure can be modified separately from the others. If the derivative payoff offsets part of the cash-flow variation associated with that factor, the combined position becomes less sensitive to it. If the payoff reinforces the existing exposure, sensitivity to the factor can increase.

Hedging is one application of this mechanism. By taking a derivative position with an offsetting payoff, a participant can reduce a selected exposure while leaving other components of the original position largely unchanged.

\subsection{Risk Sharing and Market Completeness}

Risk sharing becomes more flexible when market participants can trade claims with different state-contingent payoffs. A firm seeking more stable cash flows may transfer part of an exposure to another participant willing to bear it in exchange for compensation or because the exposure fits the participant's other positions. Derivatives facilitate such transfers by creating contractual payoffs tied to particular future states.

The extent of risk sharing depends on the range of payoffs that can be constructed from traded claims. Consider a simple economy with two future states, denoted by $u$ and $d$. Suppose two traded securities have payoff vectors

\begin{equation}
\mathbf{x}_1=
\begin{bmatrix}
x_{1u}\\
x_{1d}                                                      
\end{bmatrix},
\qquad
\mathbf{x}_2=
\begin{bmatrix}
x_{2u}\\
x_{2d}
\end{bmatrix}.
\label{eq:ch11-two-security-payoffs}
\end{equation}

A portfolio with holdings $a$ and $b$ in the two
securities produces the following state-contingent
payoff vector:

\begin{equation}
\mathbf{x}_p
=
a\mathbf{x}_1+b\mathbf{x}_2
=
\begin{bmatrix}
ax_{1u}+bx_{2u}\\
ax_{1d}+bx_{2d}
\end{bmatrix},
\label{eq:ch11-spanning-portfolio}
\end{equation}

Positive values of $a$ and $b$ represent long
positions in the respective securities, whereas
negative values represent short positions.

If the two payoff vectors are linearly independent,
they span the two-dimensional state-payoff space.
Consequently, suitable portfolio holdings can
generate any payoff vector across the two states.
A market in which every relevant state-contingent
payoff can be constructed from available traded
claims is called a \emph{complete market}.

If the available claims cannot generate all relevant state-contingent payoffs, the market is incomplete. Some desired patterns of cash flows cannot then be constructed exactly, and participants may continue to bear residual exposures after trading the available instruments.

In actual markets, transaction costs, trading restrictions, nontradable risks, model uncertainty, and a limited range of contracts may prevent some desired state-contingent payoffs from being constructed.

Derivatives can expand the range of available state-contingent payoffs by introducing claims linked to risks not fully represented by existing securities. This can increase the range of risks that can be transferred and shared across market participants. Derivatives do not create additional economic resources; they redistribute future cash flows and risks across states and participants.

\subsection{Financial Innovation and New Claims}

In an incomplete market, some desired state-contingent payoffs cannot be constructed from existing claims. Financial innovation can partly address this gap by creating new derivatives linked to particular prices, rates, indices, events, or maturities. These new claims can expand the range of exposures that can be transferred or hedged.

New claims can be designed by changing either the reference basis or the payoff rule. A contract may be linked to a new commodity index, volatility measure, or credit event. Its payoff structure may also be redesigned through caps, floors, barrier features, or path-dependent conditions. By altering what the contract is linked to or how its payoff is determined, innovation changes the state-contingent cash flows available to market participants.

A closer match, however, does not eliminate all risk. A customized claim may reduce basis risk by tracking the underlying exposure more closely, but it may trade less actively or create greater dependence on a particular counterparty. A standardized claim may offer greater liquidity while providing a less exact match to the exposure being managed. Financial innovation therefore changes not only which risks can be transferred, but also the form of the residual risks that remain.

\section{No-Arbitrage and Replication}

Derivative markets rely on two related mechanisms \cite{BlackScholes1973,Merton1973,HarrisonKreps1979}
to maintain pricing consistency. No-arbitrage
restricts the relative prices of claims according
to their state-contingent payoffs, while replication
uses portfolios of traded securities to construct
equivalent payoffs. Together, these mechanisms
provide the basis for valuing derivative contracts.

\subsection{Arbitrage Opportunities and Price Consistency}

The law of one price requires derivative positions
with identical state-contingent cash flows to have
the same current price. A meaningful comparison
requires matching underlying references, notional
amounts, maturities, payment dates, settlement
currencies, and other payoff-determining contractual
terms. Differences in these features may produce
different cash flows and do not necessarily indicate
inconsistent pricing.

An arbitrage opportunity can take two forms.
The first involves a zero initial net investment
and nonnegative future payoffs in every state,
with a strictly positive payoff in at least one
state. The second generates a positive initial
cash inflow without creating any future net payment
obligation. Both forms allow a trader to obtain
a gain without bearing an offsetting future loss.

\paragraph{First Type: Arbitrage from Payoff Dominance.}

A first form of arbitrage arises when two
derivative claims have the same current price,
but one provides at least as much as the other
in every future state and strictly more in at
least one state. Suppose

\begin{equation}
P_A=P_B,
\qquad
X_A(s)\geq X_B(s), \quad \forall s,
\label{eq:ch11-payoff-dominance}
\end{equation}
with strict inequality in at least one possible
state. An investor can purchase $A$ and sell $B$
short. The strategy requires zero initial net
investment:

\begin{equation}
\Pi_0=P_B-P_A=0.
\end{equation}

Its future net payoff satisfies

\begin{equation}
\Pi_T(s)=X_A(s)-X_B(s)\geq 0,
\qquad \forall s,
\end{equation}
with a strictly positive payoff in at least
one state. This strategy generates nonnegative
future payoffs without requiring an initial
investment and constitutes an arbitrage
opportunity.

Under the no-arbitrage condition, a claim that
provides at least as much as another in every
relevant state and strictly more in at least one
state must have a strictly higher current price,
provided the corresponding trading positions
can be established.

\paragraph{Second Type: Arbitrage between Equivalent-Payoff Claims.}

Consider two traded derivative positions, $A$ and
$B$, with identical contractual cash flows at the
same dates in every relevant future state. Let
$P_A$ and $P_B$ denote their current prices, and
let $X_A(s)$ and $X_B(s)$ denote their respective
payoffs in state $s$. Their payoff equivalence
requires

\begin{equation}
X_A(s)=X_B(s),
\qquad \forall s.
\label{eq:ch11-equivalent-derivative-payoffs}
\end{equation}

If $P_A<P_B$, an investor can purchase $A$ and
sell $B$ short. The initial net cash inflow is

\begin{equation}
\Pi_0=P_B-P_A>0,
\end{equation}

while the future net payoff is

\begin{equation}
\Pi_T(s)=X_A(s)-X_B(s)=0,
\qquad \forall s.
\end{equation}

The strategy generates a positive initial cash
inflow without creating a future net payment
obligation. If $P_A>P_B$, reversing the two
positions produces the corresponding opportunity.
Under the no-arbitrage condition, equivalent
derivative positions must therefore have the
same current price.

\begin{example}{Arbitrage from Payoff Dominance}
{ch11-payoff-dominance-arbitrage}

Consider two derivative claims with the same
maturity and settlement currency. Their current
prices and payoffs across two possible states are:

\begin{center}
\begin{tabular}{lcc}
\toprule
\textbf{Position}
& \textbf{State 1}
& \textbf{State 2} \\
\midrule
Claim A (price 9) & 12 & 20 \\
Claim B (price 9) & 10 & 20 \\
\midrule
Long A, short B & 2 & 0 \\
\bottomrule
\end{tabular}
\end{center}

Purchasing A and selling B short requires no
initial net investment. The resulting position
generates a positive payoff in state 1 and zero
payoff in state 2. The equal current prices are
therefore inconsistent with the no-arbitrage
condition. In an arbitrage-free market, claim A
must have a higher current price than claim B.

\end{example}

In actual markets, theoretical arbitrage opportunities
may not be fully exploitable. Transaction costs,
bid--ask spreads, short-sale restrictions, funding
and collateral requirements, and counterparty risk
affect the feasibility and net returns of arbitrage
strategies. 

Price consistency may also be assessed against
portfolios of traded securities when no single
traded claim provides an equivalent payoff.


\subsection{Replication and Synthetic Positions}

Replication concerns the construction of a portfolio
whose state-contingent payoffs match those of a target
financial claim. A portfolio that reproduces the
payoffs of another claim is called a
\emph{synthetic position}. When replication is exact,
the two positions have identical future cash flows
and must have the same current value under
no-arbitrage conditions.

Consider a market with $J$ traded securities and
$S$ possible future states. Let $X_j(s)$ denote the
payoff of security $j$ in state $s$, and let $P_j$
denote its current price. A portfolio holding
$\theta_j$ units of each security generates the
payoff

\begin{equation}
X_p(s)
=
\sum_{j=1}^{J}\theta_j X_j(s),
\qquad s=1,\ldots,S.
\label{eq:ch11-replicating-portfolio-payoff}
\end{equation}

Suppose a target claim has payoff $X(s)$ in state
$s$. Replication requires the portfolio to satisfy

\begin{equation}
\sum_{j=1}^{J}\theta_j X_j(s)
=
X(s),
\qquad s=1,\ldots,S.
\label{eq:ch11-replication-conditions}
\end{equation}

In matrix notation, this condition is

\begin{equation}
\mathbf{X}\boldsymbol{\theta}
=
\mathbf{x},
\label{eq:ch11-replication-matrix}
\end{equation}
where $\mathbf{X}$ is the matrix of traded-security
payoffs, $\boldsymbol{\theta}$ is the vector of
portfolio holdings, and $\mathbf{x}$ is the target
claim's payoff vector.

The current cost of the replicating portfolio is

\begin{equation}
P_p
=
\sum_{j=1}^{J}\theta_j P_j.
\label{eq:ch11-replicating-portfolio-cost}
\end{equation}

If the portfolio exactly reproduces the target
claim's contractual cash flows, no-arbitrage
requires the claim's current price $P_X$ to equal the
replication cost:

\begin{equation}
P_X=P_p
=
\sum_{j=1}^{J}\theta_j P_j.
\label{eq:ch11-replication-pricing}
\end{equation}

If the target claim trades above the replication
cost, an investor can purchase the replicating
portfolio and sell the claim. If it trades below
the replication cost, the reverse position applies.
In either case, the difference in current prices
can be captured while the future payoffs offset
in every state.

\begin{example}{Replicating a Contingent Protection Claim}
{ch11-contingent-protection-replication}

Consider an economy with three possible future
states: expansion, slowdown, and severe downturn.
Three securities are traded in the market. Their
current prices are RMB~44.20, RMB~65.00, and
RMB~48.00, respectively. Their payoffs per unit
across the three states are shown below.

\begin{center}
\begin{tabular}{@{}lccc@{}}
\toprule
\textbf{State}
& \textbf{Security 1}
& \textbf{Security 2}
& \makecell{\textbf{Risk-free}\\\textbf{security}} \\
\midrule
Expansion       & 100 & 20  & 50 \\
Slowdown        & 50  & 50  & 50 \\
Severe downturn & 20  & 100 & 50 \\
\bottomrule
\end{tabular}
\end{center}

Suppose a business-interruption protection claim
pays RMB~0 in an expansion, RMB~50 in a slowdown,
and RMB~100 in a severe downturn. Its payoff vector
is

\[
\mathbf{x}
=
\begin{bmatrix}
0\\
50\\
100
\end{bmatrix}.
\]

A portfolio of the three traded securities can
reproduce this payoff. Let $\theta_1$, $\theta_2$,
and $\theta_3$ denote the respective holdings.
The replication conditions are

\begin{align*}
100\theta_1+20\theta_2+50\theta_3&=0,\\
50\theta_1+50\theta_2+50\theta_3&=50,\\
20\theta_1+100\theta_2+50\theta_3&=100.
\end{align*}

Solving these equations gives

\[
\theta_1=-0.625,\qquad
\theta_2=0.625,\qquad
\theta_3=1.
\]

The replicating portfolio therefore consists of
a short position in 0.625 units of Security 1,
a long position in 0.625 units of Security 2,
and a long position in one unit of the risk-free
security.

Its current cost is

\begin{align*}
P_p
&=-0.625(44.20)+0.625(65)+48\\
&=61.
\end{align*}

The portfolio reproduces the protection claim's
payoffs in all three states. Under the stated
no-arbitrage conditions, the current value of
the claim is therefore RMB~61.

\end{example}

Replication is distinct from hedging in a broader
sense. A hedge may reduce the sensitivity of an
existing position to selected risk factors without
eliminating all state-dependent differences.
Exact replication, by contrast, requires the
portfolio to match the target claim's payoff in
every relevant state.

\section{Derivative Market Architecture and Risk Infrastructure}

Derivative market architecture determines how trading relationships are formed, how exposures are offset, and how counterparty risk is redistributed. Standardization and centralized trading can concentrate liquidity and improve price discovery, while collateral, netting, and clearing limit bilateral exposure and reorganize the transmission of losses across market participants.

\subsection{Trading Venues: Exchange-Traded and Over-the-Counter Markets}

Derivatives can trade on organized exchanges or through over-the-counter (OTC) arrangements. Exchange-traded derivatives generally use standardized contract terms, including contract size, maturity, settlement terms, and other specifications. Standardization allows many participants to trade contracts with common terms and makes positions more interchangeable.

OTC derivatives are negotiated between counterparties, often under standardized legal documentation while allowing economic terms to be tailored. Notional amount, maturity, reference variable, and settlement provisions can therefore be adjusted to match a particular exposure more closely. Greater customization, however, can make contracts less interchangeable and positions more difficult to transfer.

The distinction is not absolute. Some OTC derivatives use highly standardized terms and may be centrally cleared, while electronic trading platforms have reduced the traditional separation between exchange and dealer-based execution. Trading venue, contract design, and clearing arrangement are therefore related but distinct dimensions of derivative-market organization.

Major derivative exchanges provide concrete examples of how standardized contracts are organized within formal trading venues. Table~\ref{tab:ch11-major-derivative-exchanges} presents selected examples from China and major international markets.

\begin{table}[H]
\centering
\caption{Examples of Major Exchange-Traded Derivative Markets}
\label{tab:ch11-major-derivative-exchanges}

\begin{tabularx}{\textwidth}{p{3.8cm} p{3.2cm} X}
\toprule
\textbf{Exchange} &
\textbf{Market Focus} &
\textbf{Representative Derivatives} \\
\midrule

China Financial Futures Exchange (CFFEX) &
Financial derivatives &
Equity-index futures and options; government-bond futures \\

Shanghai Futures Exchange (SHFE) &
Commodity derivatives &
Metals, energy, and other commodity futures and options \\

Dalian Commodity Exchange (DCE) &
Commodity derivatives &
Agricultural, industrial, and chemical futures and options \\

Zhengzhou Commodity Exchange (ZCE) &
Commodity derivatives &
Agricultural, energy, and industrial futures and options \\

Guangzhou Futures Exchange (GFEX) &
Commodity and emerging derivatives &
Contracts linked to selected industrial and emerging-market exposures \\

Shanghai and Shenzhen Stock Exchanges &
Securities derivatives &
Exchange-traded ETF and other approved options \\

\addlinespace
CME Group &
Financial and commodity derivatives &
Interest-rate, equity-index, foreign-exchange, and commodity derivatives \\

Eurex &
Financial derivatives &
Equity-index, interest-rate, and other derivatives \\

ICE &
Financial and commodity derivatives &
Energy, interest-rate, equity, and other derivatives \\

\bottomrule
\end{tabularx}
\end{table}

\subsection{Liquidity and Price Discovery}

Derivative-market design affects where trading activity and information are concentrated. Standardized contracts can bring participants with similar exposures into common instruments, which may support trading liquidity and make market prices more observable. Customized contracts can match individual exposures more closely, but trading is often distributed across more heterogeneous positions.

Liquid markets can also contribute to price discovery. When participation is broad and transactions occur frequently, quoted and traded prices can incorporate new information and provide reference values for related instruments and underlying markets. This informational role depends on market depth, participation, and trading conditions; standardization alone does not ensure an active or informative market.


\subsection{Counterparty Exposure, Collateral, and Netting}

A derivative can acquire positive or negative market value as market conditions change. When a bilateral contract has positive value to one party, that party has a claim on the counterparty and may suffer a loss if the counterparty fails to perform. Counterparty exposure therefore changes with the market value of the contract and can evolve over time.

Counterparty credit risk is distinct from the market risk of the derivative. Market risk arises from changes in the underlying price, rate, or other risk factor, while counterparty credit risk arises from the possibility that a contractual obligation will not be fulfilled. The two risks can interact when market movements increase exposure at the same time that the counterparty's credit quality deteriorates.

Margin requirements provide a mechanism for
controlling counterparty exposure as derivative
values change. Two forms of margin serve different
purposes. Initial margin is collateral posted to
cover potential future exposure that may arise
during the period required to close out or replace
a defaulted position. Variation margin reflects
changes in the current market value of a contract
and is transferred between counterparties as
these values change.

In exchange-traded futures markets, contracts are
marked to market at the end of each trading day.
Daily gains and losses are calculated using the
change in the settlement price and credited to
or debited from the participants' margin accounts.
For a long futures position involving $n$ contracts
with contract multiplier $m$, the daily variation
margin is

\begin{equation}
VM_t=nm(F_t-F_{t-1}),
\label{eq:ch11-daily-variation-margin}
\end{equation}

where $F_t$ and $F_{t-1}$ denote the current and
previous settlement prices, respectively. A short
position has the opposite variation margin.

Daily settlement transfers gains and losses as
they arise rather than allowing market-value
changes to accumulate as unsecured contractual
exposures. Initial margin provides additional
protection against potential losses during the
close-out period following a participant's default.

Collateral and netting are two important mechanisms for reducing unsecured bilateral exposure. Collateral places assets behind an obligation as contract values change. Legally enforceable netting allows positive and negative positions between the same counterparties to be considered together rather than treated separately. Netting reduces the exposure to which collateral may need to be applied, while collateral reduces the unsecured portion of the remaining obligation.

\subsection{Central Clearing and the Redistribution of Counterparty Risk}

Collateral and netting reduce counterparty exposure within bilateral relationships. Central clearing changes the structure of those relationships. Once an eligible derivative is accepted for clearing, a central counterparty (CCP) is interposed between the original parties, and each clearing member faces the CCP for the cleared obligation. Exposures that would otherwise remain bilateral are therefore reorganized within the clearing system.

\begin{figure}[htbp]

\centering

\begin{tikzpicture}[
    >=Latex,
    dealer/.style={
        draw,
        rounded corners,
        minimum width=20mm,
        minimum height=8mm,
        align=center,
        font=\small
    },
    ccp/.style={
        draw=FinanceBlue,
        fill=FinanceBlue!8,
        rounded corners,
        minimum width=24mm,
        minimum height=12mm,
        align=center,
        font=\small
    },
    bilateral/.style={
        Latex-Latex,
        line width=0.7pt,
        draw=FinanceBlue
    },
    title/.style={
        font=\small\itshape
    }
]

% Left panel: bilateral structure
\node[title] at (2.5,3.5) {Bilateral market structure};

\node[dealer] (A1) at (0.8,2.4) {Dealer A};
\node[dealer] (B1) at (4.2,2.4) {Dealer B};
\node[dealer] (C1) at (0.8,0.4) {Dealer C};
\node[dealer] (D1) at (4.2,0.4) {Dealer D};

\draw[bilateral] (A1) -- (B1);
\draw[bilateral] (A1) -- (C1);
\draw[bilateral] (A1) -- (D1);
\draw[bilateral] (B1) -- (C1);
\draw[bilateral] (B1) -- (D1);
\draw[bilateral] (C1) -- (D1);

\node[align=center, font=\scriptsize]
at (2.5,-0.55)
{Multiple bilateral exposures};

% Arrow between structures
\draw[->, line width=0.9pt, gray!70]
(5.25,1.4) -- (6.25,1.4);

% Right panel: centrally cleared structure
\node[title] at (9.0,3.5) {Centrally cleared structure};

\node[dealer] (A2) at (7.3,2.5) {Dealer A};
\node[dealer] (B2) at (10.7,2.5) {Dealer B};
\node[dealer] (C2) at (7.3,0.3) {Dealer C};
\node[dealer] (D2) at (10.7,0.3) {Dealer D};

\node[ccp] (CCP) at (9.0,1.4) {Central\\Counterparty};

\draw[line width=0.8pt, FinanceBlue] (A2) -- (CCP);
\draw[line width=0.8pt, FinanceBlue] (B2) -- (CCP);
\draw[line width=0.8pt, FinanceBlue] (C2) -- (CCP);
\draw[line width=0.8pt, FinanceBlue] (D2) -- (CCP);

\node[align=center, font=\scriptsize]
at (9.0,-0.55)
{Exposures reorganized through the CCP};

\end{tikzpicture}

\caption{Bilateral and centrally cleared counterparty relationships}

\label{fig:ch11-bilateral-central-clearing}

\end{figure}

Central clearing can also support multilateral netting across cleared positions. Obligations that would otherwise remain in separate bilateral relationships can be offset within the clearing system, reducing some gross exposures.

Central clearing therefore redistributes counterparty risk rather than eliminating it. Counterparty relationships become organized around the CCP rather than dispersed across many bilateral pairs, making the CCP a central node in the cleared market.

%========================================================
% Summary
%========================================================

\TailModule{summary}{Summary}

Derivative markets provide an institutional framework
for trading contracts whose cash flows depend on
underlying assets, rates, indices, or specified
events. These contracts create contingent claims
that allocate payments across possible future states.
By separating selected risk exposures from broader
business and financial positions, derivatives
facilitate risk transfer and sharing among market
participants. The range of attainable contingent
claims depends on the available traded instruments
and the completeness of the market.

Prices of contingent claims are linked through
no-arbitrage relationships among traded positions.
Payoff dominance imposes restrictions on relative
prices, while claims with identical state-contingent
payoffs must have the same current price. When a
claim can be replicated by a portfolio of traded
securities, its price must equal the cost of that
portfolio. The availability of traded instruments
therefore determines which claims can be replicated
exactly and, for claims that cannot, the range of
prices consistent with no-arbitrage.

Derivative-market architecture shapes trading,
settlement, and counterparty risk management.
Exchange-traded and OTC markets differ in
standardization and contractual flexibility.
Margin requirements, collateral, and netting
manage counterparty exposure, while central
clearing reorganizes bilateral exposures through
a central counterparty. These arrangements
redistribute risk without eliminating it.

%========================================================
% Key Terms
%========================================================

\TailModule{keyterms}{Key Terms}
Contingent claims; state-contingent payoffs;
exposure decomposition; risk unbundling; risk sharing;
market completeness; financial innovation;
law of one price; arbitrage; payoff dominance;
replication; synthetic positions;
exchange-traded derivatives; over-the-counter (OTC)
derivatives; liquidity; price discovery;
counterparty exposure; margin requirements;
collateral and netting; central clearing and
central counterparty (CCP).

%========================================================
% Financial AI Lab
%========================================================

\TailModule{aiprompt}{Financial AI Lab}

%========================================================
% Additional Financial AI Lab exercises
%========================================================

%========================================================
% Financial AI Lab
%========================================================


\textbf{Lab: Audit Counterparty Exposure and Risk
Management in Derivative Markets}

Ask an AI system to assess the counterparty exposure
of a firm's OTC derivatives portfolio. Consider a
firm with the following positions, valued from the
firm's perspective on the same date. All amounts
are measured in RMB millions.

\begin{center}
\begin{tabular}{lcc}
\toprule
\textbf{Counterparty}
& \textbf{Contract}
& \textbf{Market Value} \\
\midrule
Dealer A & Contract 1 & $+6.0$ \\
Dealer A & Contract 2 & $-4.0$ \\
Dealer B & Contract 3 & $+3.0$ \\
\bottomrule
\end{tabular}
\end{center}

The two contracts with Dealer A are subject to
the same legally enforceable close-out netting
agreement. Dealer A has posted variation margin
of RMB~1.2 million to the firm, measured after
applicable collateral haircuts. The contract with
Dealer B is subject to a separate netting agreement,
and no collateral has been received.

Require the AI system to calculate the firm's
current unsecured counterparty exposure and
evaluate how different risk-management arrangements
affect this exposure.

Audit the answer in four stages.

\begin{enumerate}

\item \textbf{Calculate net counterparty exposure.}

Determine the net current exposure to each dealer
after applying legally enforceable netting.
Account for collateral received and explain why
positions with different counterparties cannot
automatically be offset against each other.

\item \textbf{Distinguish market risk from
counterparty credit risk.}

Explain how changes in underlying prices affect
the market value of the contracts and the firm's
counterparty exposure. Identify why a positive
replacement value creates exposure to the
counterparty's failure to perform.

\item \textbf{Evaluate margin and collateral
arrangements.}

Distinguish initial margin from variation margin.
Explain how daily mark-to-market and variation
margin transfers affect current exposure, while
initial margin provides protection against
potential future exposure during the close-out
period. Identify risks that may remain despite
collateralization.

\item \textbf{Assess the role of central clearing.}

Explain how interposing a central counterparty
changes the original bilateral relationships.
Discuss multilateral netting, margin requirements,
default resources, and the concentration of
counterparty risk within the clearing system.
Distinguish the redistribution of counterparty
risk from the elimination of market risk.

\end{enumerate}

The final report should verify the exposure
calculations, identify the legal and operational
assumptions underlying netting and collateral,
and distinguish risk reduction from risk transfer
and risk concentration.

%========================================================
% Discussion Questions
%========================================================

\TailModule{questions}{Discussion Questions}

\begin{enumerate}

\item How can a derivative modify a selected financial
exposure without changing the firm's underlying
business activity?

\item How does market completeness affect the ability
to replicate contingent claims, and what pricing
uncertainty may remain when exact replication is
not feasible?

\item Why can a customized OTC contract provide a
closer match to an underlying exposure while
increasing the importance of collateral, margin
requirements, and bilateral netting arrangements?

\item How does central clearing reorganize counterparty
relationships, and why does this make the central
counterparty's risk management and resilience
important?

\end{enumerate}

%========================================================
% Applied Problems
%========================================================

\TailModule{appprobs}{Applied Problems}

\begin{enumerate}

\item \textbf{State-contingent payoffs and forward hedging.}

A commodity producer expects to sell one unit of
a commodity at date $T$. The spot price at that
date may be 70, 100, or 140. The producer enters
a cash-settled short forward contract with a
delivery price of 100 and the same maturity.

\begin{enumerate}

\item Calculate the payoff of the short forward
position in each state.

\item Calculate the producer's combined cash flow
from the physical sale and the forward position
in each state.

\item Explain how the forward contract changes
the distribution of the producer's cash flows
across states.

\end{enumerate}


\item \textbf{Exposure decomposition and risk unbundling.}

A firm will purchase 500 units of an imported
commodity at date $T$. The commodity price in
foreign currency is $P_T$, and the domestic-currency
price of one unit of foreign currency is $S_T$.
The firm's domestic-currency purchase cost is

\[
C_T=500P_TS_T.
\]

\begin{enumerate}

\item Identify the two market risk factors affecting
the firm's purchase cost.

\item Explain how changes in each risk factor affect
the firm's domestic-currency cash outflow.

\item Suppose the firm uses a derivative linked
only to the commodity price $P_T$. Which exposure
is modified, and which exposure remains?

\end{enumerate}


\item \textbf{Replication of a contingent claim.}

Two securities are traded in a two-state market.
Their current prices and payoffs are as follows:

\begin{center}
\begin{tabular}{lccc}
\toprule
\textbf{Security}
& \textbf{Current Price}
& \textbf{Up State}
& \textbf{Down State} \\
\midruleSecurity A & 8 & 12 & 6 \\
Security B & 9 & 4 & 10 \\
\bottomrule
\end{tabular}
\end{center}

A contingent claim pays 20 in the up state and
2 in the down state.

\begin{enumerate}

\item Determine the holdings of securities A
and B required to replicate the contingent claim.

\item Identify which positions are long and
which are short.

\item Calculate the current cost of the
replicating portfolio.

\item Explain why the claim's current value
must equal the replication cost under the
no-arbitrage condition.

\end{enumerate}


\item \textbf{Arbitrage opportunities and price consistency.}

Consider two pairs of derivative claims.

For the first pair, claims A and B have identical
payoffs of 15 in state 1 and 8 in state 2. Their
current prices are 13 and 10, respectively.

For the second pair, claims C and D have the
same current price of 11. Their payoffs are:

\begin{center}
\begin{tabular}{lcc}
\toprule
\textbf{Claim}
& \textbf{State 1}
& \textbf{State 2} \\
\midrule
Claim C & 14 & 20 \\
Claim D & 12 & 20 \\
\bottomrule
\end{tabular}
\end{center}

\begin{enumerate}

\item For the first pair, identify the arbitrage
strategy and calculate its initial cash flow
and future net payoffs.

\item For the second pair, identify the arbitrage
strategy and calculate its initial cash flow
and future net payoffs.

\item Why does an investment with a high expected
return not necessarily constitute an arbitrage
opportunity?

\end{enumerate}


\end{enumerate}
%========================================================
% Further Reading
%========================================================

\TailModule{furtherread}{Further Reading}

\begin{itemize}

\item Duffie, Darrell, and Chi-Fu Huang (1985),
``Implementing Arrow-Debreu Equilibria by Continuous Trading of Few Long-Lived Securities,''
\emph{Econometrica}, 53(6), 1337--1356.

\end{itemize}

\ChapterTail{summary,keyterms,aiprompt,questions,appprobs,furtherread}
